\documentclass[11pt]{article}
\usepackage[margin=1in]{geometry}
\usepackage[T1]{fontenc}
\usepackage{lmodern,amsmath,amssymb,booktabs,enumitem,array,xcolor,longtable}
\usepackage[most]{tcolorbox}
\usepackage[colorlinks,linkcolor=blue!50!black]{hyperref}
\setlength{\parskip}{5pt}\setlength{\parindent}{0pt}
\newcommand{\SP}{\ensuremath{\spadesuit}}\newcommand{\DI}{\ensuremath{\diamondsuit}}\newcommand{\CL}{\ensuremath{\clubsuit}}
\newtcolorbox{say}{breakable,before upper={\setlength{\parskip}{5pt}},colback=blue!3,colframe=blue!40!black,left=6pt,right=6pt}
\newtcolorbox{act}{breakable,before upper={\setlength{\parskip}{5pt}},colback=orange!8,colframe=orange!70!black,left=6pt,right=6pt}
\newcommand{\slidehead}[3]{\subsection*{Slide #1. #2 \hfill {\normalfont\small #3}}}
\newcommand{\cue}[1]{{\small\itshape\color{black!65}[#1]}}
\title{Talk Script and Structure\\[4pt]\large Range-Vectorized Discounted Regret Minimization for No-Limit Hold'em}
\author{}\date{}
\begin{document}
\maketitle
\vspace{-3em}
This is a word-for-word script for the 60-minute talk with slides \texttt{tex/talk.pdf} and the live demonstration
on the site. Blue boxes are what to say. Orange boxes are what to do. Italic cues in brackets are stage directions.
The script is written at about 140 words per minute. Speak it in your own words once you know it; the numbers are
the parts to keep exact.

\tableofcontents

\section{Structure and timing}
\begin{center}
\begin{tabular}{rllrr}
\toprule
& part & slides & minutes & clock at end\\ \midrule
1 & The problem and how error is measured & 1--5 & 7 & 0:07\\
2 & The algorithm: CFR and DCFR & 6--7 & 7 & 0:14\\
3 & One iteration: traversal and leaves & 8--11 & 9 & 0:23\\
4 & Cost, memory, precision, parallelism & 12--15 & 9 & 0:32\\
5 & Convergence and validation & 16--18 & 6 & 0:38\\
6 & Live demonstration & 19 & 9 & 0:47\\
7 & Extensions, limitations, conclusions & 20--28 & 8 & 0:55\\
8 & Questions & & 5 & 1:00\\ \bottomrule
\end{tabular}
\end{center}

\textbf{The story in one breath.} A solver is one simple iteration repeated a few hundred times. The number of
repetitions is about the same everywhere, so everything depends on making one iteration cheap, accurate and
parallel. Parts 3 and 4 are the heart; do not let parts 1 and 2 run long.

\textbf{If you are running late.} At 0:25 instead of 0:23: shorten Memory to two sentences. At 0:40 instead of
0:38: use the five-minute demonstration. In part 7, slides 23 and 24 can each be cut to one sentence.

\textbf{Before you start.} Add your name to the title slide. Open the site in a second window with the Postflop page
loaded and the K72 example selected; open the Preflop page in another tab so the equity tables are already
downloaded. Close other heavy programs.

% =====================================================================================
\section{Part 1. The problem and how error is measured (7 minutes)}

\slidehead{1}{Title}{1 min}
\begin{say}
Thank you all for coming. My talk is about poker solvers, the programs that compute how poker should be played.

I want to start with a claim. A poker solver is one numerical method, called discounted counterfactual regret
minimization, applied to one structured optimization problem. The method itself fits on a single slide. So the
interesting question is not what the algorithm is. The question is what makes it fast and accurate. How is one step
organized? What does it cost? How much precision does it need? And how do you split it across processors?

Those are numerical analysis questions, and they are what this talk is about. I built a solver to have something
real to measure, and you will see it run live later.
\end{say}

\slidehead{2}{Background and scope}{1.5 min}
\begin{say}
Some background. In 2017 two programs beat professional players at two-player no-limit hold'em, and in 2019 a
program beat professionals at the six-player game. The engine inside all of them is a family of algorithms called
counterfactual regret minimization, or CFR.

Since then, commercial solvers have become standard study tools for players. You give them a situation: the cards
on the board, the hands each player might hold, and the bet sizes you allow. They return the equilibrium strategy.

What I built is a solver of that kind that runs in a web browser. It handles the flop, the turn and the river. It
has the main features of the commercial tools. And it uses one algorithm throughout, discounted CFR.

On the right is its output: one player's strategy, for every starting hand, in one situation. Green is check and
red is bet. A square that is part green and part red is a hand that mixes. The whole talk is about how numbers like
these are computed.
\end{say}

\slidehead{3}{Outline}{0.5 min}
\begin{say}
Here is the plan. The top row is the core: the problem, the algorithm, how one iteration is organized, what it
costs, precision, parallelism and convergence. Then I will switch to the browser for a demonstration, and finish
with the extensions built on the same core and what comes next.
\end{say}

\slidehead{4}{The postflop subgame}{2 min}
\begin{say}
Let me set up the problem, and I will assume no poker knowledge.

Each player holds two private cards. Public cards, the board, are dealt in stages: three cards called the flop,
then one called the turn, then one called the river, with betting after each stage.

A solver works on a subgame. You specify the board, the money in the pot, the money left to bet, and two more
things. First, a range for each player. You do not know your opponent's hand, but you know which hands are
plausible, and a range is simply a weight for every possible hand. It is a vector. Second, a bet tree: the few bet
sizes we allow, which makes the game finite.

How long are these vectors? With three board cards there are at most 1176 two-card hands, and 1081 on the river.

The unknowns are an action distribution for every hand at every decision point. My running example is the most
common situation in online poker, on the board king, seven, two. The table shows the sizes. A river problem has
seventeen decision points. The flop problem has fifty-seven thousand, each with about four hundred and seventy
hands. That is tens of millions of unknowns, because every possible turn card and every possible river card opens
its own copy of the later betting.
\end{say}

\slidehead{5}{Zero-sum games, equilibrium and exploitability}{2 min}
\begin{say}
What does it mean to solve this? Two-player poker is zero-sum: what I win, you lose. If I describe my strategy by a
vector x and yours by y, my expected winnings are x transpose A y, linear in each.

Von Neumann's minimax theorem says there is a pair of strategies where neither player can improve alone. That is a
Nash equilibrium, and it is what a solver computes. For the optimizers in the room, this is a saddle-point problem.

Now the most important definition of the talk. An iterative method never reaches the exact equilibrium, so we need
a measure of error. It is called exploitability. Take the current strategies. Ask how much I could gain by
switching to my best possible reply, and how much you could gain by switching to yours. Average the two. It is
zero exactly at equilibrium, and we report it as a percentage of the pot.

Two things make it a good error measure. It has a meaning: it is money a perfect opponent could take from you. And
it is computable without knowing the answer. In practice, 0.3 percent of the pot counts as solved.

One caveat: if the casino takes a fee from the pot, the game is not zero-sum and the theorems I show do not apply.
The solver still runs and reports the same quantity.
\end{say}
\cue{Clock check: 0:07.}

% =====================================================================================
\section{Part 2. The algorithm (7 minutes)}

\slidehead{6}{Regret matching and counterfactual regret minimization}{3.5 min}
\begin{say}
Here is the algorithm. It rests on one idea: regret.

After each round of play, for each action, ask: how much better would I have done with this action than with what I
actually did? That is the regret. Add it up over time. Then use a very simple rule, the formula on the right: play
each action in proportion to its positive accumulated regret. If I keep wishing I had bet more, I bet more.

A quick example with rock, paper, scissors. I start uniform and you play rock. Paper would have won, so paper has
regret plus one, and next round I play paper. You play scissors, and now rock has the most regret. The strategies
chase each other. What converges is the average over time.

Poker has many decision points, and CFR applies this rule separately at every decision point and for every hand.
The word counterfactual refers to the payoff it uses: the value of an action, weighted by the probability that the
opponent and the cards actually lead to this point.

The theorem says this works. The exploitability of the average strategies is bounded by a constant over the square
root of the number of iterations T.

Look at the constant. It contains the number of information sets, which for our flop problem is twenty-seven
million. If you put in the numbers, after ten thousand iterations the bound is around a hundred million chips, for a
pot of five and a half chips. So the theorem guarantees convergence and says nothing useful about speed. That gap
between theory and practice is why I will measure convergence later.
\end{say}

\slidehead{7}{Discounted CFR}{3.5 min}
\begin{say}
The version I use is discounted CFR, from Brown and Sandholm in 2019. The motivation is simple. Early iterations
are bad, because the strategies start out random. Plain CFR remembers those early regrets forever.

Discounted CFR forgets. Before adding the new regret, it multiplies the old sum by a factor a little below one.
Positive regrets are multiplied by t to the one and a half over t to the one and a half plus one. Negative regrets
are halved. And in the average strategy, iteration t gets weight t squared. So after ten iterations the first
iteration counts for a quarter of a percent, where a plain average would give it ten percent.

A nice feature is that plain CFR and the well-known variant CFR-plus are special cases of the same formula. In my
code they differ only in three parameters, which will make the comparison fair.

I also update the players in alternation: the second player responds to the first player's new strategy. If you
know iterative methods for linear systems, this is Gauss--Seidel in place of Jacobi.

The last line on this slide is there because of a mistake I made, and I think it is the most instructive moment of
the project. If the opponent never reaches some part of the tree, all values there are zero. So I skipped those
parts. It seemed harmless. But skipping also skipped the discounting there, so I was no longer running the
algorithm on the slide. On a small test game the error after a thousand iterations was two hundred and fifty times
larger than it should have been. And it still converged, slowly, so nothing looked broken.

The fix is to visit those parts, apply the discounts, and skip only the leaf computations. And there is now a test
that the numbers match a version with no shortcut, bit for bit. The lesson: an optimization is only correct if it
leaves the iterates unchanged, and that needs a test.
\end{say}
\cue{Clock check: 0:14.}

% =====================================================================================
\section{Part 3. One iteration: traversal and leaves (9 minutes)}

\slidehead{8}{Range-vectorized traversal of the public tree}{2.5 min}
\begin{say}
Now the heart of the talk: how one iteration is organized.

The solver works on what is called the public tree. It contains only the things both players see: the bets and the
board cards. The private hands are not in the tree. Instead, at every node, every quantity is a vector with one
entry per hand.

One iteration is one walk through this tree for each player. On the way down, I carry the opponent's reach vector:
for each of his hands, the probability that he plays it this way. On the way up, I return a value vector: what each
of my hands is worth.

The box lists what happens at each node. At the opponent's decision, I multiply his reach, entry by entry, by his
strategy for each action and go down. At my own decision, I go down every action, combine the values with my
strategy, and update my regrets. At a chance node I average over the cards. At a leaf I compute payoffs.

Why does this matter? A naive program walks the tree once for every pair of hands, more than two hundred thousand
walks. This version walks it twice per iteration, and at each node it runs a short loop over about five hundred
numbers stored side by side in memory, which is what processors are best at. The cost of a decision node is
proportional to the number of hands times the number of actions.

And notice: the products A times y and A transpose times x from the earlier slide are exactly what these walks
compute. The matrix A is never formed.
\end{say}

\slidehead{9}{Chance nodes}{1.5 min}
\begin{say}
A short slide on chance nodes, where a card is dealt. The value of my hand is the average over the cards that can
come. The question is the weight.

Fix both players' hands. On the flop, three board cards and four private cards are known. That leaves 45 cards,
each equally likely. In general K is 48 minus the number of board cards.

The point of the lemma is that this number is the same for every pair of hands. So a chance node is a plain sum
times one over K, with no bookkeeping per hand. It is easy to get wrong: using one over 49, the number of cards not
on the board, makes every value about eight percent too small. I check this against exact solutions later.
\end{say}

\slidehead{10}{Fold leaves in linear time}{2.5 min}
\begin{say}
Now the leaves, where the hand ends. This is where the main trick is.

There is a complication called card removal. Two players cannot hold the same card. If I hold the ace of spades,
you do not. So when I compute the value of my hand, I may only count your hands that share no card with mine.

The obvious way is a double loop: for each of my hands, go through all of yours and skip the clashes. That is n
squared operations per leaf. With a thousand hands that is a million operations, at tens of thousands of leaves, in
every iteration.

Here is the better way for a fold, where I only need the total weight of your compatible hands. Take your total
weight, tau. Subtract the weight of your hands that contain my first card. Subtract the weight of those that contain
my second card. One hand, the one identical to mine, has now been subtracted twice, so add it back. That is
inclusion and exclusion.

To use it I need your total and one sum for each of the 52 cards. That is one pass over your hands, then one pass
over mine. Linear, not quadratic.

At the bottom is the linear algebra view. The table of which hands are compatible is the all-ones matrix, minus one
small correction per card, plus a correction for identical hands. A simple matrix plus a few corrections can be
applied quickly. Exploiting structure in place of storing a matrix is the recurring theme here.
\end{say}

\slidehead{11}{Showdown leaves in linear time}{2.5 min}
\begin{say}
A showdown is harder, because the payoff depends on who has the better hand. The idea is to sort.

Sort your hands by strength, once for each board. Now the total weight of your hands that are weaker than mine is a
running total along the sorted list. So I sweep up both sorted lists once, keeping a running total, and each of my
hands reads off the weight it beats and the weight it ties. Card removal is handled inside the sweep with the same
52 card sums as before.

The formula says: assume I lose to everything compatible, then correct for the hands I beat and the hands I tie.

The table shows it works. These are microseconds per leaf. Look at the last column, the double loop: when n
doubles, it roughly quadruples. The sweep, in the middle, roughly doubles. At a thousand and eighty-one hands the
sweep takes fourteen microseconds and the double loop takes four thousand: three hundred times faster. And the two
agree to twelve digits.

One flop iteration evaluates more than forty thousand showdowns per player, so this is the difference between two
seconds and several minutes per iteration.
\end{say}
\cue{Clock check: 0:23.}

% =====================================================================================
\section{Part 4. Cost, memory, precision, parallelism (9 minutes)}

\slidehead{12}{A cost model for one iteration}{2.5 min}
\begin{say}
So what does an iteration cost overall? I want a model I can test, not just a big-O statement.

The model has two terms. The first counts the entries of the regret table: hands times actions, summed over all
decision points. The second counts the hands at the showdowns that were evaluated. Each has a constant in
nanoseconds.

I timed 252 iterations on eight different trees and fitted the two constants by least squares: thirty nanoseconds
per table entry and seventeen per hand at a showdown.

A fit always looks good on its own data, so the last column is the honest test. I leave one tree out, fit on the
other seven, and predict the one left out. On the two large flop trees the error is seven and ten percent. On the
smallest trees it is up to thirty percent, where an iteration takes a few milliseconds and timing noise is large.
The model also under-predicts trees dominated by leaves, so some per-leaf overhead is missing from it.

What you should take from this slide: time per iteration is proportional to the size of the tree, about two to
three seconds for these flop trees on one core. And adding a bet size multiplies the tree, because the new branch
carries its own copy of everything below it.
\end{say}

\slidehead{13}{Memory}{1.5 min}
\begin{say}
Memory follows the same count. Each entry stores two numbers, a regret and a strategy sum, at four bytes each. Eight
bytes times the number of entries. The simplest flop tree needs 431 megabytes. Add raises and sizes and you pass
three gigabytes, which is more than a browser tab gets. Memory, not time, is what limits tree size in a browser.

Two practical points. The numbers live in a few very large arrays, and each node holds a slice. Allocating one
array per node took thirty-two seconds for a large tree; this takes under one. And the interface predicts the
memory before building, by building one card's subtree and multiplying. Where I could build the full tree to
check, the estimate was right to a tenth of a percent.
\end{say}

\slidehead{14}{Floating-point precision}{2.5 min}
\begin{say}
Next, precision. I store everything in single precision, about seven significant digits, to halve the memory. Is
that safe?

The risk with single precision is adding a small number to a large running sum: the small number gets rounded
away. Here the sums grow slowly, thanks to the discounting, and a new term would only be lost after tens of
millions of iterations. We run a few hundred. So the theory says it is safe.

The plot is the experiment. I ran all four combinations: regrets in single or double precision, strategy sums in
single or double. All four curves reach the same exploitability, down to two ten-thousandths of a percent of the
pot. That is a thousand times below the practical target. Storage precision is not what limits accuracy.

There is one subtle finding. The exploitabilities agree, but the strategies from single and double precision differ
a little, by a bit under one percent on average. Is that rounding damage? I ran a control. Two double-precision
runs, where I changed the input ranges in the seventh decimal place. They differ by just as much.

So it is not about storage. Any tiny perturbation sends the method to a slightly different strategy that is
equally good. In numerical analysis terms, the strategy is ill-conditioned with respect to exploitability. Picture a
wide flat valley floor: every point is almost equally low, but the points are far apart. The solver guarantees how
low you are, not where on the floor you stand.
\end{say}

\slidehead{15}{Parallel partition by the first chance card}{2.5 min}
\begin{say}
Last in this part: parallelism. The tree has a natural place to cut. After the first betting round a card is
dealt, and what happens below one card is independent of what happens below another.

So I give each processor a share of the cards. A coordinator keeps the small top of the tree. One pass is three
steps. The coordinator walks the top and records the reach vectors at each point where a card is dealt. Each worker
computes the values under its own cards and returns one sum. The coordinator divides by K and finishes.

Very little data moves, a few kilobytes of vectors, against millions of operations per card.

The model is simple: the time is the coordinator's time plus the number of cards of the busiest worker times the
time per card. On the flop there are 49 turn cards, so with two workers the busier one has 25.

The table is measured in the browser. On the flop: about 2.3 seconds on one core, 1.2 seconds with two workers. A
speedup of 1.9 on two cores, and the model predicts it to within the timing noise. My test machine has two cores, so
I cannot show more; the model predicts about a third of a second with eight, and I label that an extrapolation.

One honest detail. The parallel result is not bit-for-bit identical to the sequential one, because the card sums
are added in a different order. They agree to fourteen digits.
\end{say}
\cue{Clock check: 0:32.}

% =====================================================================================
\section{Part 5. Convergence and validation (6 minutes)}

\slidehead{16}{Convergence of CFR, CFR$^+$ and DCFR}{2 min}
\begin{say}
We know what one iteration costs. How many do we need?

These plots show exploitability against iterations, on logarithmic axes, for a river, a turn and a flop problem.
Blue is plain CFR, orange is CFR-plus, green is discounted CFR. The dashed line is the 0.3 percent target.

On a log-log plot a straight line of slope minus k means the error behaves like one over T to the k. The theorem
promised slope minus one half. What we measure at the end of these runs is about minus one for CFR, and between
minus 1.4 and minus 2.1 for discounted CFR. On the flop after 150 iterations, CFR is at 1.6 percent of the pot and
discounted CFR is at 0.05 percent.

Two cautions. These are slopes over a finite window, not proven rates. And nobody can currently explain why the
method beats its own worst-case bound by this much. That is an open problem.
\end{say}

\slidehead{17}{Iterations to a fixed accuracy}{2 min}
\begin{say}
One problem per street is thin evidence, so here are thirteen: five river, five turn, three flop. The table gives
the median number of iterations to reach each accuracy, with the range.

To reach 0.3 percent: about 95 on the river, 170 on the turn, 71 on the flop. Between different boards the count
varies by less than a factor of two.

The lower table fixes one flop and changes the tree. The first two trees differ by a factor of a hundred in size and
need the same forty iterations. What increases the count is more betting options, not more nodes.

This is the conclusion that organizes the whole talk. The number of iterations is roughly constant, a hundred or
two. The cost per iteration ranges over three orders of magnitude from river to flop. So the cost per iteration
decides the running time, and that is why I spent so long on it.
\end{say}

\slidehead{18}{Validation}{2 min}
\begin{say}
Before the demonstration: how do I know any of this is correct?

Small zero-sum games can be solved exactly by a completely different method, linear programming. So I built small
games a second time, separately, in Python, solved them with a linear programming solver, and compared the value of
the game.

On a turn game, the two values agree to seven digits. On a flop game with two levels of dealt cards, to six. On a
larger turn game with 40 and 35 randomly chosen hands, to eight digits. These tests exercise the chance weighting,
the card removal and the showdowns together.

I compare values and not strategies because of the flat valley: strategies are not unique, the value is.

The other rows check components. The parallel solver against the sequential one. The shortcut against no shortcut,
bit for bit. The fast leaves against the double loop.

What is missing is a comparison with an established solver on a full-size flop tree. I have not done that yet, and
I will come back to it under limitations.
\end{say}
\cue{Clock check: 0:38. Switch to the browser.}

% =====================================================================================
\section{Part 6. Live demonstration (9 minutes)}

\slidehead{19}{Demonstration}{9 min}
\begin{say}
Let me show you all of this running. Everything you will see runs in this browser tab; nothing is sent to a server.
\end{say}

\subsection*{Step 1. Solve a flop (about 2.5 minutes)}
\begin{act}
Show the Postflop page with the K72 example. Point at the board, the two range grids and the bet tree. Point at the
estimate line under the tree. Press \textbf{Build and solve}.
\end{act}
\begin{say}
This is the running example. The board is king, seven, two. On the left are the two ranges, and below them the bet
tree with one bet size per street.

Before I build anything, this line gives the estimate from the memory slide: about fifty-seven thousand decision
points and 431 megabytes.

I press build and solve. \cue{press} It is now building the tree in the worker threads. You can see the number of
workers here.

And here is the convergence plot from the talk, live. Each point costs two best-response computations, so I only
evaluate every ten iterations. The axis is logarithmic and the dashed line is the 0.3 percent target. Each
iteration is one or two seconds on this laptop, and we need about seventy of them.
\end{say}
\cue{While it runs, about one minute: keep talking.}
\begin{say}
While it runs: every iteration is exactly the two tree walks I described. The workers are each handling their share
of the 49 turn cards and sending back sums. The reason this takes a minute or two and a commercial solver takes
seconds is that this is plain JavaScript. They use compiled code and exploit suit symmetry.
\end{say}
\cue{If it has not reached the target after about 90 seconds, press \textbf{Stop} and say: ``I will stop it here at
about [value] percent, which is enough to browse.''}

\subsection*{Step 2. Browse the solution (about 2 minutes)}
\begin{act}
The strategy browser opens. Read the action buttons. Click \textbf{check}. In the card picker click the
5\SP. Then switch the view to \textbf{EV}, then \textbf{Equity}, then back. Point at one square.
\end{act}
\begin{say}
The solve has stopped, and this is the strategy browser. In this tree the first decision is the button's, and the
buttons show how often he checks and how often he bets across his whole range.

I click check. Now a card is dealt, so I choose one: the five of spades. \cue{click} Every turn card has its own
solution; that is the copy of the tree under each card.

This grid is the big blind's strategy on that turn. Each square is a starting hand. Green is check, red is bet, and
the split inside a square is the mixing frequency. Strong hands and some weak hands bet; the middle mostly checks.
That is the bluffing structure game theory predicts, and nobody programmed it in.

I can switch to expected value, in chips, or to equity. And if I point at one square, the panel on the right lists
each individual combination with its own strategy. The grid is a summary. The solver works with the individual
hands, which is where card removal shows up.
\end{say}

\subsection*{Step 3. Lock a node (about 2 minutes)}
\begin{act}
Open \textbf{Node locking}. Press \textbf{Edit and lock}. Choose the bet action, set the frequency to 100\%, and
drag across one row of medium hands. Press \textbf{Lock this strategy}, then \textbf{Continue}. When it stops,
click the bet action to show the button's response.
\end{act}
\begin{say}
Now a feature players use constantly. Suppose I know my opponent makes a specific mistake. Here, let me make the big
blind bet these medium hands every time. I paint them, and lock.

Locking freezes this strategy. When I press continue, the solver re-solves everything else around it.

Watch the exploitability: it jumps, because the old solution no longer fits, and comes back down. \cue{wait} And
now look at the button's response to the bet: it has changed. Typically he continues more often, because the
betting range now contains medium hands he can attack.

Mathematically, the solver is now finding the equilibrium of a restricted game. Nothing else in the code changes:
locked regrets are simply not updated.
\end{say}

\subsection*{Step 4. Preflop (about 2 minutes)}
\begin{act}
Open the \textbf{Preflop and multiway} tab. The preset is heads-up push or fold, 10 bb. Press \textbf{Build and
solve}. Point at the first action frequencies. Then choose \textbf{Six-max, 20 bb, with antes}, solve, and show the
\textbf{First-in ranges} table.
\end{act}
\begin{say}
The same algorithm also runs before the flop. Here strategies are per class of starting hand, 169 of them.

This is the simplest preflop game: each player has ten big blinds and can only go all-in or fold. \cue{press} It
converges in about a second. The small blind goes all-in with 58 percent of hands, and the big blind calls with 37
percent. Those are the published equilibrium numbers, and a linear program gives the same values to five digits.

Now six players. \cue{select and press} With more than two players there is no convergence guarantee, so the page
reports how much the players could still gain by deviating. This table shows how often each seat enters the pot. It
is a sensible strategy, but it is not a proven equilibrium, and I say so on the page.
\end{say}
\cue{Switch back to the slides. Clock check: 0:47.}

\subsection*{The five-minute version}
Do step 1 (stop the solve at 60 seconds), step 2 without the EV and equity views, and only the push-or-fold solve of
step 4. Skip locking and say: ``The next slide shows locking with numbers.''

\subsection*{If something goes wrong}
\begin{itemize}[nosep]
\item \textbf{The page will not load or the solve fails.} Go back to the slides. Slide 19 has a screenshot of the
browser; describe step 2 from it and move on.
\item \textbf{The solve is very slow.} Stop it early. Any exploitability under about 2\% gives a browsable solution.
\item \textbf{Out of memory.} Reload and choose the turn example spot (K72 with a 5); it needs little memory and
solves in seconds.
\end{itemize}

% =====================================================================================
\section{Part 7. Extensions, limitations, conclusions (8 minutes)}

\slidehead{20}{Node locking}{1.25 min}
\begin{say}
Here is locking with numbers, on a river problem. In equilibrium the button never bets a third of the pot after a
check. I lock him to always make that bet.

Against the lock, the big blind folds 29 percent of the time in place of 49, and her result improves by ten chips a
hand. Always betting is an expensive mistake, and the solver quantifies it.

One subtlety. The 49 percent is a response to a bet that never happens in equilibrium, and game theory does not pin
down behaviour in situations that never arise. And the last row: the locked game is no harder to solve.
\end{say}

\slidehead{21}{Aggregate flop reports}{1.25 min}
\begin{say}
Players also ask questions across all flops, such as how often to bet on average. There are 22,100 flops, but
suits are interchangeable, which leaves 1,755 that are really different.

A report solves a sample of them and takes a weighted average. This is numerical integration: a quadrature rule
over flops. The sample is stratified by the type of board, and that gives the error bars in the table. With 25, 49
or 95 flops the answers agree within those error bars. The middle column is the error from not solving each flop
exactly, and it is smaller than the sampling error.
\end{say}

\slidehead{22}{Preflop: card removal between hand classes}{1.5 min}
\begin{say}
Before the flop, as you saw, the vectors have length 169. Card removal now acts between classes. If I hold ace-king,
it is half as likely that you hold two aces. The matrix M records these fractions exactly. It is all ones except
where two classes share a rank, about a quarter of the entries, so multiplying by it is again a sum with a sparse
correction.

The table is the check against linear programming for the all-in-or-fold game at four stack sizes. The values agree
to five digits.

The limitation is that a hand reaching the flop is settled as if both players were all-in. Betting after the flop is
not modelled here.
\end{say}

\slidehead{23}{Payoff errors and exploitability}{1 min}
\begin{say}
The preflop solver depends on a table of winning probabilities that is estimated by random sampling, so it has
small errors. We are solving a slightly wrong game. How wrong is the solution?

The lemma answers in the worst case: if every payoff is off by at most delta, the solution is off by at most
delta in exploitability. This is classic error propagation.

Then I measured it, with two independent tables. The actual effect is ten thousand to a hundred thousand times
smaller than the bound, because random errors cancel when you average over a range of hands.
\end{say}

\slidehead{24}{Three-way showdowns: a low-rank factorization}{0.75 min}
\begin{say}
With three players all-in, the probabilities form a table with 169 cubed entries, and using it directly is too slow.
Each slice is a symmetric matrix, so I keep its sixteen most important eigen-components. That is five times cheaper,
with an error of about a quarter of a percent in the quantity we need.
\end{say}

\slidehead{25}{Multiway convergence}{0.5 min}
\begin{say}
With more than two players there is no theorem. So I measure. The all-in-or-fold games converge essentially
exactly. A realistic six-player game stalls at a few hundredths of a big blind. Useful, but not an equilibrium.
\end{say}

\slidehead{26}{Limitations}{0.5 min}
\begin{say}
The limitations, plainly. It is slow next to commercial solvers. The fast convergence has no theory. The preflop
model ignores later betting. Parallel scaling is measured on two cores only. And the flop solver has not yet been
compared with an established solver on a full-size tree; my exact checks go up to 40 hands per player.
\end{say}

\slidehead{27}{Next steps}{0.5 min}
\begin{say}
So the first next step is that comparison. After that: the high-precision finish with a Newton-type method that I
presented last time, which is where new research would go. Then engineering: suit symmetry, compiled code, and
pruning done in a way that keeps the iterates exact.
\end{say}

\slidehead{28}{Conclusions}{0.75 min}
\begin{say}
To conclude. One method, discounted CFR, solves all of these games.

What makes it fast is structure. Vectors over hands turn an iteration into simple loops. Inclusion and exclusion
and a sorted sweep make every leaf linear in place of quadratic.

What makes it trustworthy is measurement. A cost model tested on held-out trees, a precision experiment with a
control, a parallel version checked against the sequential one, and exact solutions to compare with.

And the number of iterations is about the same everywhere, so the cost of one iteration is what decides everything.

The site, the code and every experiment are at the address on the slide. Thank you. I am happy to take questions.
\end{say}
\cue{Clock check: 0:55.}

% =====================================================================================
\section{Part 8. Questions you should expect}
\begin{description}[leftmargin=1em,itemsep=4pt]
\item[Why not just solve it as a linear program?] ``For small games I do, as a check. For the flop the matrix would
have a dense five-hundred by five-hundred block at each of about seventy thousand leaves. It cannot be stored. CFR
only needs products with the matrix, and the tree walk computes those without forming it.''
\item[Why does DCFR beat its bound by so much?] ``Nobody knows. The bound is a worst case over all games. Poker
trees clearly have more structure. It is an open problem and I list it under next steps.''
\item[Is the last strategy or the average the answer?] ``The average. The current strategy keeps moving; the
weighted average is what converges and what the theorem is about.''
\item[How does this compare with PioSolver in speed?] ``Much slower: seconds per flop iteration against a small
fraction of a second. They use compiled, vectorized code and solve only one of each group of equivalent cards. The
algorithmic ideas, the vectors and the linear-time leaves, are the same.''
\item[Why trust it without a comparison to a commercial solver?] ``The exact linear programming checks cover every
component together on games up to 40 hands per player, and the exploitability I report is computed independently of
the iteration. A full-size comparison is the first thing I will do next, and I say on the slide that it is missing.''
\item[If strategies are not unique, what use is the output?] ``Every output is within the stated exploitability, so
any of them is nearly unbeatable. What is not determined is fine detail, such as whether a hand bets 30 or 35
percent. The value of the game and the broad structure are stable.''
\item[Why single precision if double is safer?] ``Half the memory, and memory is the limit in a browser. The
experiment shows the same exploitability with either.''
\item[Why is the parallel result not identical?] ``Floating-point addition depends on the order of the terms in the
last digits. The workers add card values in a different order. The difference is in the fourteenth digit.''
\item[What happened with the skipped subtrees exactly?] ``The values there are zero, so skipping the values is fine.
But DCFR also discounts the stored regrets every iteration, everywhere. Skipping the visit skipped the discount, so
those regrets were stale when the opponent started reaching that part again.''
\item[Does it handle more than two players after the flop?] ``No. Postflop is two players. Preflop handles up to
six, without a convergence guarantee.''
\item[What does the Newton method add?] ``DCFR gets three or four digits cheaply and then slows down. A Newton-type
method started from that point reached thirteen digits on river problems in my earlier work. Extending it to the
flop is the research direction.''
\end{description}

\newpage
\section{One-page cheat sheet of numbers}
\begin{center}\small
\begin{tabular}{ll}
\toprule
Hands avoiding the board & 1176 flop, 1128 turn, 1081 river\\
Running example & K\SP 7\DI 2\CL, 488 and 460 hands, pot 5.5, stacks 97.5\\
Decision points & 17 river, 1{,}349 turn, 56{,}842 flop\\
Target accuracy & 0.3\% of the pot\\
DCFR parameters & $(1.5,\ 0,\ 2)$\\
Bug factor & 250 times worse (0.64 against 0.0025 chips)\\
Cards left at a chance node & $K=48-|\text{board}|$: 45 flop, 44 turn\\
Leaf at 1081 hands & sweep 13.7 $\mu$s, double loop 4272 $\mu$s, about 300$\times$\\
Cost model & 30 ns per entry, 17 ns per showdown hand; held-out error 7 and 10\%\\
Memory & 8 bytes per entry; 431 MB for the simplest flop tree\\
Precision & same exploitability down to $2\times10^{-4}$\%; strategies differ by about 0.007\\
Parallel, flop & 2317 ms sequential, 1212 ms with 2 workers, speedup 1.9\\
Slopes & CFR about $-1$; DCFR $-1.4$ to $-2.1$; theory $-0.5$\\
Iterations to 0.3\% & river 95, turn 170, flop 71 (medians)\\
LP checks & agree to 6--9 digits; largest game 40 and 35 hands\\
Push or fold, 10 bb & shove 58\%, call 37\%\\
Flops & 22{,}100 total, 1{,}755 distinct\\
Lock example & fold 49\% $\to$ 29\%, value $-16.2\to-6.2$ chips\\ \bottomrule
\end{tabular}
\end{center}
\end{document}
